\subsection{Inputs}

\footnotesize
\begin{longtable}{@{}P{1.5cm}P{2.2cm}P{2.2cm}P{3.7cm}P{1.2cm}P{4.2cm}@{}}
\caption{Inputs of the model.}\label{tab:terms}\\
\toprule
\textbf{Symbol} & \textbf{Name} & \textbf{Type} & \textbf{Represents} & \textbf{Directly observed} & \textbf{Obtained how} \\
\midrule
\endfirsthead
\toprule
\textbf{Symbol} & \textbf{Name} & \textbf{Type} & \textbf{Represents} & \textbf{Directly observed} & \textbf{Obtained how} \\
\midrule
\endhead
\bottomrule
\endfoot
\mm{$x$}{Assets-to-liabilities ratio}{A real number above one.}{Total assets over total liabilities, the state of the model (A1).}{Yes}{Recorded as total assets over total liabilities in the Call Report, \path{ASSET} over \path{LIAB} (parent map G1).}
\mm{$\pi$}{Risky exposure}{A function of the state with values between zero and the cap over the state.}{The share of assets in the risky asset (A1, A2). No row derives it; M11 takes it at the cap and M12 bounds every compliant choice.}{Yes}{Recorded as risk-weighted assets over total assets, \path{RWAJT} over \path{ASSET}, under risk weights of zero on the riskless and one on the risky asset (parent map G2). On the panel its median is 0.69 against a cap of 2.38 at the median state, so the cap does not bind there.}
\mm{$\sigma$}{Asset volatility}{A positive real number.}{The volatility of the risky asset, and the letter of the diffusion of the state (A1, A3).}{No}{Assumed at the baseline of \citet{bayraktar2026}, 0.08 (parent map G8). The volatility of the book ratio is computed on the panel and is 4.1 times smaller than the model's at the recorded exposure, so it is not this object.}
\mm{$\sigma_L$}{Liability volatility}{A positive real number.}{The volatility of the shock to liabilities (A1).}{No}{Assumed at the baseline of \citet{bayraktar2026}, 0.03 (parent map G9).}
\mm{$c$}{Return-liability correlation}{A real number of absolute value below one.}{The correlation between the risky return and the liability shock (A1).}{No}{Assumed at the baseline of \citet{bayraktar2026}, 0.20 (parent map G9).}
\mm{$q$}{Required capital}{A positive real number.}{The capital regulation requires, as a ratio to liabilities, the zero of the capital buffer (A1).}{No}{Assumed. A floor on the leverage ratio gives a candidate, 3/97 for the 3\% Tier 1 leverage ratio when Tier 1 equals book equity, and the bundle does not fix it (parent map G7). Buffer variance does not depend on it.}
\mm{$a_1$}{Solvency floor}{A real number between zero and the liquidity haircut.}{The floor on equity over risky assets, the slope of the solvency branch of the regulatory cap (A2).}{Yes}{Recorded in the Basel standard as the CET1 minimum of 4.5\% of risk-weighted assets, the baseline of \citet{bayraktar2026}; which tier it stands for is not fixed (parent map G4).}
\mm{$a_2$}{Run-off rate}{A real number below one.}{Net 30-day outflows as a share of liabilities, the intercept of the liquidity branch (A2).}{No}{Assumed at 0.05, the standard's run-off factor for stable retail deposits. The standard sets a factor per class of liability, so one rate for all liabilities is the model's reduction (parent map G5).}
\mm{$a_3$}{Liquidity haircut}{A real number between the solvency floor and one.}{The haircut on the risky asset in high-quality liquid assets, the slope of the liquidity branch (A2).}{No}{Assumed at 0.30, the baseline of \citet{bayraktar2026}, which is none of the standard's haircuts (parent map G6).}
\mm{$\rho_L$}{Liability-net discount rate}{A positive real number.}{The discount rate of the problem written per unit of liabilities (A4).}{No}{Assumed from the shareholders' discount rate, 0.12 at the baseline of \citet{bayraktar2026} (parent map G10).}
\mm{$\lambda_S$}{Asymmetry parameter}{A real number at least one.}{How much more the objective weighs a shortfall below the reference than a surplus above it (A4). At one the shortfall penalty vanishes and the model has issuance costs alone (M8).}{No}{Never observed. It enters the rows through the shortfall penalty only. M3 identifies it locally with the fixed issuance cost from the dividend barrier and the recapitalisation trigger, a statement within the model, since on the panel neither threshold is measured as one (the rows of the dividend barrier and the recapitalisation trigger below).}
\mm{$K$}{Fixed issuance cost}{A positive real number.}{The cost of an issue that does not scale with its size, per unit of liabilities (A5).}{No}{Never observed. Its counterpart is the part of an issue's cost that does not scale with the amount raised, which \citet{peura2006} set from implied issue cost and frequency rather than record (L8). M3 identifies it locally with the asymmetry parameter from the two thresholds.}
\mm{$\kappa$}{Proportional issuance cost}{A non-negative real number.}{The cost of an issue per unit raised (A5).}{No}{Assumed at the baseline of \citet{bayraktar2026}, 0.01 for dilution and 0.02 per unit issued (parent map G11).}
\end{longtable}
\normalsize

\subsection{Internal objects}

\footnotesize
\begin{longtable}{@{}P{1.5cm}P{2.2cm}P{2.2cm}P{3.7cm}P{1.2cm}P{4.2cm}@{}}
\caption{Internal objects of the model.}\\
\toprule
\textbf{Symbol} & \textbf{Name} & \textbf{Type} & \textbf{Represents} & \textbf{Directly observed} & \textbf{Obtained how} \\
\midrule
\endfirsthead
\toprule
\textbf{Symbol} & \textbf{Name} & \textbf{Type} & \textbf{Represents} & \textbf{Directly observed} & \textbf{Obtained how} \\
\midrule
\endhead
\bottomrule
\endfoot
\mm{$z$}{Capital buffer}{A real number, $z=\log((x-1)/q)$.}{The log of equity over required capital, zero where equity equals required capital and falling without limit towards the distress boundary (A1, M9).}{No}{Computed from the assets-to-liabilities ratio and required capital by its formula, whose inverse is \lean{CapGeometry.zmap_inverse} and which is zero where equity equals required capital, \lean{Coordinate.zmap_at_requirement}. The parent bundle records its discrete-time form as Tier 1 capital over 8.5\% of risk-weighted assets (parent map G18).}
\mm{$R$}{Reference}{A real number, $R=1+q$.}{The state at which equity equals required capital, below which the shortfall penalty applies (A4). The capital buffer is zero there.}{No}{Computed from required capital by its formula. The capital buffer vanishes at it, \lean{Coordinate.zmap_at_requirement}; the penalty is zero at or above it, \lean{QVI_Part1.shortfall_zero_at_or_above_reference}, and positive below it for an asymmetry parameter above one, \lean{QVI_Part1.shortfall_pos_below_reference}. Its counterpart is the required capital of the solvency floor's row, the level at which prompt corrective action's restrictions begin (parent ledger A9).}
\mm{$u$}{Regulatory cap}{A function of the state, $u(x)=\min\{(x-1)/a_1,(x-a_2)/a_3\}$.}{The largest risky exposure regulation allows at each state, the smaller of its two branches (A2).}{No}{Computed from the assets-to-liabilities ratio and the cap geometry by its formula. The solvency branch is the smaller up to the crossing point, \lean{CapGeometry.solvency_binds}, and the liquidity branch beyond it, \lean{CapGeometry.liquidity_binds}.}
\mm{$\bar x$}{Crossing point}{A real number above one, $\bar x-1=a_1(1-a_2)/(a_3-a_1)$.}{The state at which the two branches of the regulatory cap are equal (M10).}{No}{Computed from the cap geometry by its formula, \lean{CapGeometry.xbar_sub_one}.}
\mm{$\Lambda$}{Shortfall penalty}{A function of the state, $\Lambda(x)=(\lambda_S-1)(R-x)^+$.}{The term of the objective that charges a shortfall below the reference (A4). Zero when the asymmetry parameter is one, \lean{QVI_Part1.lambda_asymmetry_vanishes_at_one}.}{No}{Never observed. It is a term of the objective, and whose cost it stands for is a reading, the shortfall-penalty paragraph of the manuscript's primitives appendix.}
\mm{$V$}{Value function}{A function of the state, twice continuously differentiable on the inaction region (A6).}{The value of the bank under its best policy, per unit of liabilities.}{No}{Never observed. Its counterpart would be the market value of equity per unit of liabilities, which the panel does not carry; \citet{daihuangkeppo2019} fit their parameters to such values (L21).}
\mm{$\mathcal MV$}{Impulse operator}{A function of the state, the supremum over targets of the value after an issue net of both issuance costs.}{The value of issuing at once to the best target (A6, M5, M7).}{No}{A device. It exists only inside the quasi-variational inequality and the proofs of M5 and M7.}
\mm{$W$}{Expected discounted shortfall}{A function of the state, $W=-\partial_{\lambda_S}V$.}{The expected discounted shortfall below the reference under the optimal policy (A6).}{No}{A device. It enters M1 to M3 through the envelope identity only.}
\mm{$N$}{Discounted injection count}{A function of the state, $N=-\partial_KV$.}{The expected discounted number of issues under the optimal policy (A6).}{No}{A device. It enters M1 to M3 through the envelope identity only; its sample counterpart, a discounted count of recorded issues, is used by no row.}
\end{longtable}
\normalsize

\subsection{Outputs}

\footnotesize
\begin{longtable}{@{}P{1.5cm}P{2.2cm}P{2.2cm}P{3.7cm}P{1.2cm}P{4.2cm}@{}}
\caption{Outputs of the model.}\\
\toprule
\textbf{Symbol} & \textbf{Name} & \textbf{Type} & \textbf{Represents} & \textbf{Directly observed} & \textbf{Obtained how} \\
\midrule
\endfirsthead
\toprule
\textbf{Symbol} & \textbf{Name} & \textbf{Type} & \textbf{Represents} & \textbf{Directly observed} & \textbf{Obtained how} \\
\midrule
\endhead
\bottomrule
\endfoot
\mm{$y^*$}{Dividend barrier}{A real number above one.}{The state at which dividends are paid (A6). M1 gives its response to the asymmetry parameter and the fixed issuance cost.}{Yes}{Recorded as the state at which a dividend is paid, from dividend quarters (\path{EQCDIV}) and the recorded state. On the panel, 56\% of bank-quarters below the solved barrier pay a dividend and payment falls as the state rises, so no barrier is measured there, and under the conservation buffer the recorded payouts reflect a rule as well as the barrier (parent map G14).}
\mm{$x_L$}{Recapitalisation trigger}{A real number above one and below the dividend barrier.}{The state at which capital is raised (A6). M2 gives its response to the two parameters.}{Yes}{Recorded as the state just before a recapitalisation, from net stock issuance (\path{EQCSXQ}) and capital from the parent holding company (\path{EQCBHCTR}). On the panel, issues start at a mean state of 1.11 against a solved trigger of 1.03 (parent map G15).}
\mm{$y_{\text{post}}$}{Injection target}{A real number between the trigger and the barrier.}{The state an issue reaches (A6). M7 says it does not depend on where the issue starts.}{Yes}{Recorded as the state just after a recapitalisation. On the panel the post-issue state rises with the pre-issue state, against the single target of M7 (parent map G15).}
\mm{$\zeta$}{Buffer variance}{A function of the state, $z'(x)^2$ times the diffusion $\sigma^2(x,\pi(x))$.}{The variance per unit of time of the capital buffer (A3).}{No}{Computed from the capital buffer, the asset volatility, the liability volatility, the return-liability correlation and the risky exposure by its formula, \lean{CapGeometry.zvol2_eq}. The parent bundle did not compute it, since a level needs the asset volatility, which is assumed (parent map G16).}
\mm{$\nu_1^2$}{Deficit limit}{A real number, $\nu_1^2=\sigma^2/a_1^2-2c\sigma\sigma_L/a_1+\sigma_L^2$.}{The limit of buffer variance at the distress boundary where the exposure binds the solvency branch (M10, M11).}{No}{Computed from the asset volatility, the liability volatility, the return-liability correlation and the solvency floor by its formula, \lean{CapGeometry.deficit_limit}. It applies where the cap binds, which on the panel it does not (the row of the risky exposure).}
\mm{$\nu_3^2$}{Surplus limit}{A real number, $\nu_3^2=\sigma^2/a_3^2-2c\sigma\sigma_L/a_3+\sigma_L^2$.}{The limit of buffer variance deep in surplus where the exposure binds the liquidity branch (M11), which needs borrowing (C3).}{No}{Computed from the asset volatility, the liability volatility, the return-liability correlation and the liquidity haircut by its formula, \lean{CapGeometry.surplus_limit}. It applies in the leveraged case only (C3).}
\end{longtable}
\normalsize
